% ============================================
% MEM-19
% Gerado de molambudos.md — build_miolo.py
% ============================================
\section*{\textbf{MEM-19} --- O Visitante}
\markboth{MEM-19 --- O Visitante}{}

\noindent \textit{Alguém veio me ver}

\noindent \textbf{1950 --- Alguém veio me ver}

\vspace{0.5em}\noindent\rule{\textwidth}{0.5pt}\vspace{0.5em}

\noindent Ninguém nunca vinha me ver.

\noindent Não tinha família. Não tinha amigos. O velho Anselmo tinha morrido no curral. Minha mãe tinha desaparecido no caminhão. O resto --- vizinhos, conhecidos --- tinha morrido na seca ou se espalhado pelo mundo.

\noindent Eu era um homem sem passado.

\noindent Por isso, quando o enfermeiro disse \textit{\textquotedblleft{}Visita para você\textquotedblright{}}, eu achei que era engano.

\noindent \textit{\textquotedblleft{}Quem?\textquotedblright{}}

\noindent \textit{\textquotedblleft{}Não sei. Uma mulher. Está na sala de visitas.\textquotedblright{}}

\noindent Uma mulher.

\noindent Minha mãe? Impossível. Ela teria mais de 70 anos. E estava morta --- ou pior.

\noindent Mas eu fui.

\noindent Atravessei o corredor. Os outros pacientes olhavam. Visita era raro no Colônia. As pessoas não visitavam loucos. Loucos eram esquecidos.

\noindent A sala de visitas era uma sala com duas cadeiras e uma mesa. Numa das cadeiras, uma mulher estava sentada.

\noindent Era idosa. Cabelos brancos, pele enrugada, mãos grossas de trabalho. Ela usava um vestido simples, azul, com um xale nos ombros.

\noindent Ela levantou os olhos quando eu entrei.

\noindent Os olhos dela eram castanhos. Cansados. Mas vivos.

\noindent \textit{\textquotedblleft{}Joaquim?\textquotedblright{}}

\noindent A voz era trêmula.

\noindent \textit{\textquotedblleft{}...Quem é?\textquotedblright{}}

\noindent Ela sorriu. Um sorriso triste.

\noindent \textit{\textquotedblleft{}Sou Maria. Maria do Rosário.\textquotedblright{}}

\noindent Minha mãe.

\noindent Eu sentei.

\noindent Fiquei olhando para ela. Tentando encontrar a mulher que me levou pela mão na estrada de Senador Pompeu. A mulher que tossia no curral. A mulher que foi levada no caminhão.

\noindent \textit{\textquotedblleft{}A senhora...\textquotedblright{}\textit{ Minha voz falhou. }\textquotedblleft{}A senhora está viva.\textquotedblright{}}

\noindent \textit{\textquotedblleft{}Estou.\textquotedblright{}}

\noindent \textit{\textquotedblleft{}Onde a senhora estava?\textquotedblright{}}

\noindent Ela baixou os olhos.

\noindent \textit{\textquotedblleft{}Fui levada para um bordel em Fortaleza. Fiquei lá 3 anos. Depois fugi. Casei de novo. Tive outros filhos. Morei em São Paulo.\textquotedblright{}}

\noindent \textit{\textquotedblleft{}E nunca...\textquotedblright{}\textit{ Eu engoli. }\textquotedblleft{}Nunca voltou?\textquotedblright{}}

\noindent \textit{\textquotedblleft{}Voltei.\textquotedblright{}\textit{ Ela me olhou. }\textquotedblleft{}Voltei para te buscar. Em 1919. Disseram que tu tinha morrido.\textquotedblright{}}

\noindent \textit{\textquotedblleft{}Mentira.\textquotedblright{}}

\noindent \textit{\textquotedblleft{}Eu sei. Descobri ano passado. Uma vizinha disse que o filho dela trabalhou aqui e que tinha um paciente do Ceará. 1.260. O mesmo nome. Vim ver.\textquotedblright{}}

\noindent \textit{\textquotedblleft{}A senhora veio me buscar?\textquotedblright{}}

\noindent Ela não respondeu.

\noindent Olhou para as mãos.

\noindent \textit{\textquotedblleft{}Joaquim, eu...\textquotedblright{}\textit{ Ela parou. }\textquotedblleft{}Eu tenho outra família. Outros filhos. Eles não sabem de ti. Não sabem que eu tive um filho antes.\textquotedblright{}}

\noindent \textit{\textquotedblleft{}E a senhora vai contar?\textquotedblright{}}

\noindent Ela não respondeu.

\noindent Eu entendi.

\noindent Ela não veio me buscar.
Ela veio se \textit{despedir}.

\noindent \textit{\textquotedblleft{}Tá\textquotedblright{}\textit{, eu disse. }\textquotedblleft{}Tá bom.\textquotedblright{}}

\noindent \textit{\textquotedblleft{}Joaquim...\textquotedblright{}}

\noindent \textit{\textquotedblleft{}Tá bom, mãe.\textquotedblright{}}

\noindent A palavra \textit{mãe} saiu sem eu pensar. Ela ouviu. Os olhos dela encheram de lágrima.

\noindent \textit{\textquotedblleft{}Me perdoa\textquotedblright{}}, ela disse.

\noindent \textit{\textquotedblleft{}Perdoar o quê?\textquotedblright{}}

\noindent \textit{\textquotedblleft{}Por ter deixado tu.\textquotedblright{}}

\noindent \textit{\textquotedblleft{}A senhora não deixou. Levaram.\textquotedblright{}}

\noindent \textit{\textquotedblleft{}Mas eu não voltei.\textquotedblright{}}

\noindent \textit{\textquotedblleft{}A senhora voltou agora.\textquotedblright{}}

\noindent Ela chorou.

\noindent Eu não chorei. A criatura dentro de mim não deixou.

\noindent Ela ficou mais 20 minutos. Contou sobre a vida dela. Os outros filhos. O marido. A casa em São Paulo. Eu ouvi. Assenti. Sorri nos lugares certos.

\noindent Quando ela foi embora, ela me abraçou.

\noindent O abraço dela era quente. Macio. Cheirava a sabão.

\noindent \textit{\textquotedblleft{}Eu vou escrever\textquotedblright{}}, ela disse.

\noindent \textit{\textquotedblleft{}Tá.\textquotedblright{}}

\noindent Ela não escreveu.

\noindent Nunca mais.

\noindent Mas, naquela noite, a criatura ficou quieta. Ela respeitou o silêncio.

\noindent \textit{\textquotedblleft{}Tu tá triste?\textquotedblright{}} ela perguntou.

\noindent \textit{\textquotedblleft{}Não.\textquotedblright{}}

\noindent \textit{\textquotedblleft{}Tá.\textquotedblright{}}

\noindent \textit{\textquotedblleft{}Ela não volta.\textquotedblright{}}

\noindent \textit{\textquotedblleft{}Eu sei.\textquotedblright{}}

\noindent \textit{\textquotedblleft{}Mas por que ela veio?\textquotedblright{}}

\noindent Pensei.

\noindent \textit{\textquotedblleft{}Para eu saber que ela existiu. Para eu saber que não fui um sonho.\textquotedblright{}}

\noindent A criatura não respondeu.

\noindent Naquela noite, pela primeira vez em anos, eu não senti fome.

\noindent Senti saudade.

\noindent Saudade de uma mãe que eu mal conhecia.
Saudade de uma vida que eu nunca tive.
Saudade de ser filho de alguém.

\noindent Saudade de ser gente.

\vspace{0.5em}
\noindent\textit{\textbf{↪ Rotas:} \rota{MEM-15} \rotasep \rota{MEM-20}}